% Unit 097 terjemahan Bahasa Indonesia dari chapter7.tex baris 1806--1967.
% Sumber: Methods of Algebra, Volume 2, commit
% 9a5803ff77dd3257484cb177f851a73770a59dd3 (CC BY 4.0).
% stable-unit-id: o014.aljabr2.chapter7.morita-theory-b
% Cakupan: Bagian 7.7, dari contoh perubahan gelanggang hingga tepat sebelum
% Lema~\ref{prop:fg-module-cat}; melanjutkan Unit 096 dan dilanjutkan Unit 098.

\hypertarget{o014.aljabr2.chapter7.morita-theory-b}{ }
% segment-id: o014.aljabr2.chapter7.morita-theory-b.e001
\begin{example}[Perubahan gelanggang]\label{eg:comonad-change-ring}
	\index[sym1]{FBA@$\mathcal{F}_{B \mid A}$}
	Tetapkan homomorfisme aljabar-$\Bbbk$ $f:A\to B$. Dalam kerangka di atas,
	ambillah $P:=B$, yang dipandang sebagai bimodul-$(A,B)$ melalui $f$. Dalam
	hal ini, adjoin kanan $\Hom_{\cated{Mod}B}(B,\mathord\cdot)$ dari
	$(\mathord\cdot)\dotimes{A}B$ isomorfik, melalui
	$\varphi\mapsto\varphi(1)$, dengan funktor pelupa
	$\mathcal{F}_{B|A}:\cated{Mod}B\to\cated{Mod}A$, dan
% segment-id: o014.aljabr2.chapter7.morita-theory-b.q001
	\[\begin{array}{|c|c|} \hline
		\text{morfisme unit }\; \eta_N & \text{morfisme kounit }\; \epsilon_M \\ \hline
		N \to \mathcal{F}_{B|A}(N \dotimes{A} B) & \mathcal{F}_{B|A}(M) \dotimes{A} B \to M \\
		x \mapsto x \otimes 1 & y \otimes b \mapsto yb \\ \hline
	\end{array} \quad \begin{array}{c}
		N: \;\text{modul-$A$ kanan,} \\
		M: \;\text{modul-$B$ kanan.}
	\end{array}\]

	Dalam pembahasan perubahan gelanggang, lazimnya $\mathcal{F}_{B|A}$
	dihilangkan dari notasi untuk menyederhanakannya. Dengan demikian, monad
	$T$ dan komonad $L$ masing-masing dapat dituliskan sebagai
% segment-id: o014.aljabr2.chapter7.morita-theory-b.q002
	\[\begin{array}{|c|c|c|} \hline
		T & \identity \to T & T^2 \to T \\ \hline
		N \mapsto N \dotimes{A} B & N \to N \dotimes{A} B & (N \dotimes{A} B) \dotimes{A} B \to N \dotimes{A} B \\
		& x \mapsto x \otimes 1 & (x \otimes b) \otimes b' \mapsto x \otimes bb' \\ \hline\hline
		L & L \to \identity & L \to L^2 \\ \hline
		M \mapsto M \dotimes{A} B & M \dotimes{A} B \to M & M \dotimes{A} B \to (M \dotimes{A} B) \dotimes{A} B \\
		& y \otimes b \mapsto yb & y \otimes b \mapsto (y \otimes 1) \otimes b \\ \hline
	\end{array}\]
	Verifikasi yang diperlukan bersifat rutin.
\end{example}

% segment-id: o014.aljabr2.chapter7.morita-theory-b.p001
Selanjutnya, tinjau suatu keadaan yang sedikit lebih umum. Pertama, untuk
sembarang modul-$B$ kanan $P$ dan $X$, pandang $B$ sebagai bimodul-$(B,B)$
untuk mendefinisikan modul-$B$ kiri
% segment-id: o014.aljabr2.chapter7.morita-theory-b.q003
\begin{equation}\label{eqn:P-vee}
	P^\vee := \Hom_{\cated{Mod}B}(P, B)
\end{equation}
serta homomorfisme modul-$\Bbbk$ kanonik
% segment-id: o014.aljabr2.chapter7.morita-theory-b.q004
\begin{equation}\label{eqn:dual-Hom-X}\begin{aligned}
	X \dotimes{B} P^\vee & \to \Hom_{\cated{Mod}B}(P, X) \\
	x \otimes \lambda & \mapsto x \lambda(\cdot).
\end{aligned}\end{equation}
\index[sym1]{Pvee@$P^\vee$}

% segment-id: o014.aljabr2.chapter7.morita-theory-b.p002
Jika $P$ adalah bimodul-$(A,B)$, maka $P^\vee$ mempunyai struktur
bimodul-$(B,A)$, $\Hom_{\cated{Mod}B}(P,X)$ mempunyai struktur modul-$A$
kanan, dan \eqref{eqn:dual-Hom-X} merupakan homomorfisme modul-$A$ kanan.

% segment-id: o014.aljabr2.chapter7.morita-theory-b.p003
Konstruksi yang sama tentu memiliki versi dengan kiri dan kanan dipertukarkan.
Karena itu, $P^{\vee\vee}$ terdefinisi, dan terdapat pula homomorfisme
kanonik modul-$B$ kanan $P\to P^{\vee\vee}$.

% segment-id: o014.aljabr2.chapter7.morita-theory-b.l001
\begin{lemma}\label{prop:AB-comonad-aux}
	Misalkan bimodul-$(A,B)$ $P$, sebagai modul-$B$ kanan, projektif dan
	dibangkitkan secara hingga. Maka $P^\vee$, sebagai modul-$B$ kiri, juga
	projektif dan dibangkitkan secara hingga, dan dalam keadaan ini:
	\begin{itemize}
		\item untuk setiap $X$, \eqref{eqn:dual-Hom-X} memberikan isomorfisme
		kanonik modul-$A$ kanan;
		\item $P\to P^{\vee\vee}$ merupakan isomorfisme.
	\end{itemize}
\end{lemma}
% segment-id: o014.aljabr2.chapter7.morita-theory-b.pf001
\begin{proof}
	Sebagai modul-$B$ kanan, tuliskan $P$ sebagai suku langsung dari
	$B^{\oplus n}$ untuk suatu $n\in\Z_{\geq0}$. Maka $P^\vee$ secara alami
	merupakan suku langsung dari $B^{\oplus n}$. Jelas bahwa penggantian $P$
	dengan $B^{\oplus n}$ dalam \eqref{eqn:dual-Hom-X} menghasilkan
	isomorfisme modul-$\Bbbk$. Karena itu, pada suku langsung $P$ kita tetap
	memperoleh isomorfisme modul-$\Bbbk$, dan dengan demikian isomorfisme
	modul-$A$ kanan. Pernyataan $P\rightiso P^{\vee\vee}$ dibuktikan dengan
	cara yang sama, yakni cukup dengan memeriksanya untuk $B^{\oplus n}$.
\end{proof}

% segment-id: o014.aljabr2.chapter7.morita-theory-b.p004
Dengan demikian, di bawah hipotesis di atas, pasangan adjoin
\eqref{eqn:AB-adjunction} dapat ditulis ulang sebagai
% segment-id: o014.aljabr2.chapter7.morita-theory-b.q005
\begin{equation}\label{eqn:AB-adjunction-2}\begin{tikzcd}
	(\cdot) \dotimes{A} P : \cated{Mod}A \arrow[shift left, r] & \cated{Mod}B: (\cdot) \dotimes{B} P^\vee . \arrow[shift left, l]
\end{tikzcd}\end{equation}

% segment-id: o014.aljabr2.chapter7.morita-theory-b.d001
\begin{definition}\label{def:bimodule-ev-coev}
	\index[sym1]{ev-coev@$\mathrm{ev}, \mathrm{coev}$}
	Misalkan bimodul-$(A,B)$ $P$, sebagai modul-$B$ kanan, projektif dan
	dibangkitkan secara hingga. Definisikan homomorfisme bimodul-$(B,B)$
% segment-id: o014.aljabr2.chapter7.morita-theory-b.q006
	\[\begin{tikzcd}[row sep=tiny]
		\mathrm{ev}: & P^\vee \dotimes{A} P \arrow[r] & B \\
		& \lambda \otimes p \arrow[mapsto, r] & \lambda(p)
	\end{tikzcd}\]
	dan homomorfisme bimodul-$(A,A)$
% segment-id: o014.aljabr2.chapter7.morita-theory-b.q007
	\[\begin{tikzcd}[row sep=tiny]
		\mathrm{coev}: & A \arrow[r] & P \dotimes{B} P^\vee \arrow[r, "{\text{Lema \ref{prop:AB-comonad-aux}}}" inner sep=0.6em, "\sim"'] & \End_{\cated{Mod}B}(P) \\
		& a \arrow[mapsto, rr] & & {[p \mapsto ap]} .
	\end{tikzcd}\]
\end{definition}

% segment-id: o014.aljabr2.chapter7.morita-theory-b.p005
Mudah diperiksa bahwa keduanya terdefinisi dengan baik. Dengan menggabungkan
\eqref{eqn:bimod-adjunction} dan Lema~\ref{prop:AB-comonad-aux}, untuk setiap
modul-$A$ kanan $N$ dan modul-$B$ kanan $M$, morfisme unit $\eta_N$ dan
morfisme kounit $\epsilon_M$ dari pasangan adjoin
\eqref{eqn:AB-adjunction-2} memenuhi
% segment-id: o014.aljabr2.chapter7.morita-theory-b.q008
\[\begin{tikzcd}[row sep=tiny]
	N \arrow[r, "\eta_N"] & \Hom_{\cated{Mod}B}(P, N \dotimes{A} P) & N \dotimes{A} P \dotimes{B} P^\vee \arrow[l, "\sim"'] \\
	x \arrow[mapsto, r] & {[p \mapsto x \otimes p]} & x \otimes \mathrm{coev}(1) \arrow[mapsto, l] \\
	M \dotimes{B} P^\vee \dotimes{A} P \arrow[r, "\sim"] & \Hom_{\cated{Mod}B}(P, M) \dotimes{A} P \arrow[r, "\epsilon_M"] & M \\
	x \otimes \lambda \otimes p \arrow[mapsto, r] & x \lambda(\cdot) \otimes p \arrow[mapsto, r] & x\lambda(p) \arrow[equal, d] \\
	& & (\identity_M \otimes \mathrm{ev})(x \otimes \lambda \otimes p).
\end{tikzcd}\]
Karena itu, kita dapat melakukan identifikasi
% segment-id: o014.aljabr2.chapter7.morita-theory-b.q009
\begin{align*}
	\eta_N = \identity_N \otimes \mathrm{coev}: \; & N \to N \dotimes{A} P \dotimes{B} P^\vee , \\
	\epsilon_M = \identity_M \otimes \mathrm{ev}: \; & M \dotimes{B} P^\vee \dotimes{A} P \to M.
\end{align*}

% segment-id: o014.aljabr2.chapter7.morita-theory-b.p006
Konstruksi di atas berlangsung pada tingkat bimodul $P$ dan $P^\vee$, dan
sama sekali tidak bergantung pada $M$ ataupun $N$.

% segment-id: o014.aljabr2.chapter7.morita-theory-b.d002
\begin{definition}\label{def:cogebroid}
	\index{modul} \index{komodul}
	Untuk sembarang aljabar-$\Bbbk$ $A$, kategori
	$(A,A)\dcate{Mod}$ merupakan kategori monoidal terhadap
	$\otimes:=\dotimes{A}$. Oleh karena itu, \S\ref{sec:alg-in-monoidal-cat}
	dan \S\ref{sec:bialgebra} mendefinisikan apa yang dimaksud dengan aljabar
	$E$ (atau koaljabar $C$) dalam $(A,A)\dcate{Mod}$, beserta modul-$E$ kiri
	dan kanan (atau komodul-$C$ kiri dan kanan). Karena definisi
	modul (atau komodul) hanya melibatkan $\otimes$ pada satu sisi, diperlukan
	perumuman kecil berikut.
	\begin{itemize}
% segment-id: o014.aljabr2.chapter7.morita-theory-b.i001
		\item Modul-$E$ kiri berarti data berikut: suatu
		$M\in\Obj(A\dcate{Mod})$ beserta morfisme dalam $A\dcate{Mod}$
		$\mu_M:E\otimes M\to M$, sedemikian sehingga diagram dalam
		Definisi~\ref{def:module-algebra} komutatif di $A\dcate{Mod}$.

% segment-id: o014.aljabr2.chapter7.morita-theory-b.i002
		\item Komodul-$C$ kiri berarti data berikut: suatu
		$M\in\Obj(A\dcate{Mod})$ beserta morfisme dalam $A\dcate{Mod}$
		$\rho_M:M\to C\otimes M$, sedemikian sehingga diagram dalam
		Definisi~\ref{def:comodule-cogebra} komutatif di $A\dcate{Mod}$.
	\end{itemize}
	Modul-$E$ kanan dan komodul-$C$ kanan didefinisikan secara serupa, dan
	morfisme antarmodul atau antarkomodul didefinisikan dengan cara standar.
\end{definition}

% segment-id: o014.aljabr2.chapter7.morita-theory-b.p007
Dengan menggunakan rangkaian hasil terdahulu untuk mendeskripsikan monad dan
komonad yang ditentukan oleh pasangan adjoin \eqref{eqn:AB-adjunction-2}, kita
langsung memperoleh hasil berikut.

% segment-id: o014.aljabr2.chapter7.morita-theory-b.pr001
\begin{proposition}\label{prop:Morita-comonad}
	Misalkan bimodul-$(A,B)$ $P$, sebagai modul-$B$ kanan, projektif dan
	dibangkitkan secara hingga, dan definisikan bimodul-$(B,A)$ $P^\vee$ seperti
	di atas. Maka:
	\begin{itemize}
		\item $P\dotimes{B}P^\vee$ merupakan aljabar dalam kategori monoidal
		$(A,A)\dcate{Mod}$. Perkaliannya diberikan oleh
% segment-id: o014.aljabr2.chapter7.morita-theory-b.q010
		\[ \identity_P \otimes \mathrm{ev} \otimes \identity_{P^\vee}: P \dotimes{B} P^\vee \dotimes{A} P \dotimes{B} P^\vee \to P \dotimes{B} P^\vee \]
		dan satuannya berasal dari
		$\mathrm{coev}:A\to P\dotimes{B}P^\vee$;
		\item $P^\vee\dotimes{A}P$ merupakan koaljabar dalam kategori monoidal
		$(B,B)\dcate{Mod}$. Komultiplikasinya diberikan oleh
% segment-id: o014.aljabr2.chapter7.morita-theory-b.q011
		\[ \identity_{P^\vee} \otimes \mathrm{coev} \otimes \identity_P: P^\vee \dotimes{A} P \to P^\vee \dotimes{A} P \dotimes{B} P^\vee \dotimes{A} P \]
		dan kounitnya berasal dari
		$\mathrm{ev}:P^\vee\dotimes{A}P\to B$.
	\end{itemize}

	Selanjutnya, tinjau monad $T$ pada $\cated{Mod}A$ yang ditentukan oleh
	pasangan adjoin \eqref{eqn:AB-adjunction-2}, serta komonad $L$ pada
	$\cated{Mod}B$ yang ditentukannya. Gunakan Definisi~\ref{def:cogebroid}
	untuk membahas modul dan komodul.
	\begin{itemize}
		\item Menentukan modul-$T$ setara dengan menentukan
		modul-$A$ kanan $N$ beserta homomorfisme yang memenuhi sifat-sifat
		standar seperti asosiativitas dan unitalitas,
% segment-id: o014.aljabr2.chapter7.morita-theory-b.q012
		\[ N \dotimes{A} (P \dotimes{B} P^\vee) \to N; \]
		dengan kata lain, setara dengan menjadikan $N$ suatu
		modul-$(P\dotimes{B}P^\vee)$.
		\item Menentukan komodul-$L$ setara dengan menentukan
		modul-$B$ kanan $M$ beserta homomorfisme yang memenuhi sifat-sifat
		standar seperti koasosiativitas dan kounitalitas,
% segment-id: o014.aljabr2.chapter7.morita-theory-b.q013
		\[ M \to M \dotimes{B} (P^\vee \dotimes{A} P); \]
		dengan kata lain, setara dengan menjadikan $M$ suatu
		komodul-$(P^\vee\dotimes{A}P)$.
	\end{itemize}
\end{proposition}
% segment-id: o014.aljabr2.chapter7.morita-theory-b.pf002
\begin{proof}
	Buktinya lebih mudah daripada pernyataannya.
\end{proof}

% segment-id: o014.aljabr2.chapter7.morita-theory-b.p008
Perhatikan bahwa, sebagai gelanggang, $P\dotimes{B}P^\vee$ tidak lain adalah
$\End_{B\dcate{Mod}}(P)$. Sebaliknya, struktur koaljabar pada
$P^\vee\dotimes{A}P$ tampak canggung, tetapi dalam penerapan sering kali
lebih nyaman. Meskipun demikian, keunggulan komodul tidak bersifat mutlak:
misalnya, agar kategori komodul menjadi kategori abelian, koaljabarnya harus
datar. Berikut hanya diberikan garis besar bukti sederhana; latihan di akhir
bab memberikan tambahan.

% segment-id: o014.aljabr2.chapter7.morita-theory-b.pr002
\begin{proposition}\label{prop:Comod-abelian}
	Misalkan $C$ koaljabar dalam $(B,B)\dcate{Mod}$. Tuliskan kategori
	komodul-$C$ kanan (atau kiri) sebagai $\cated{Comod}C$ (atau
	$C\dcate{Comod}$), dan tuliskan funktor pelupa darinya ke
	$\cated{Mod}B$ (atau $B\dcate{Mod}$) sebagai $U$. Jika $C$, sebagai
	modul-$B$ kiri (atau kanan), datar \cite[Definisi 6.9.4]{Li1}, maka
	$\cated{Comod}C$ (atau $C\dcate{Comod}$) merupakan kategori abelian,
	dan $U$ eksak setia.
\end{proposition}
% segment-id: o014.aljabr2.chapter7.morita-theory-b.pf003
\begin{proof}
	Cukup membahas versi $\cated{Comod}C$. Tinjau suatu homomorfisme komodul-$C$
	kanan $f:M\to N$, dan tuliskan kokernelnya pada tingkat modul-$B$ sebagai
	$\Coker(Uf)$. Maka, dalam $\cated{Mod}B$, terdapat diagram komutatif dengan
	baris-baris eksak
% segment-id: o014.aljabr2.chapter7.morita-theory-b.q014
	\[\begin{tikzcd}
		M \arrow[r, "Uf"] \arrow[d] & N \arrow[d] \arrow[r] & \Coker(Uf) \arrow[dashed, d] \arrow[r] & 0 \\
		M \dotimes{B} C \arrow[r, "{Uf \otimes \identity}"'] & N \dotimes{B} C \arrow[r] & \Coker(Uf) \dotimes{B} C \arrow[r] & 0,
	\end{tikzcd}\]
	dengan panah putus-putus yang ditentukan secara tunggal oleh
	funktorialitas kokernel. Mudah dibuktikan bahwa panah putus-putus tersebut
	memberi $\Coker(Uf)$ struktur komodul; objek ini kita tuliskan sebagai
	$\Coker(f)$, dan merupakan kokernel dari $f$. Pembaca dapat memeriksa
	rinciannya bila diperlukan.

	Kita ingin mengangkat $\Ker(Uf)$ menjadi komodul dengan cara yang sama.
	Kuncinya ialah memastikan bahwa
% segment-id: o014.aljabr2.chapter7.morita-theory-b.q015
	\[ 0 \to \Ker(Uf) \dotimes{B} C \to M \dotimes{B} C \xrightarrow{Uf \otimes \identity} N \dotimes{B} C \]
	eksak dalam $\cated{Mod}B$. Di sinilah diperlukan kedataran $C$ sebagai
	modul-$B$ kiri; verifikasi rutin yang tersisa serupa dengan kasus kokernel.

	Untuk membuktikan bahwa $\cated{Comod}C$ merupakan kategori abelian, masih
	perlu ditunjukkan bahwa morfisme kanonik
% segment-id: o014.aljabr2.chapter7.morita-theory-b.q016
	\[ \Coim(f) = \Coker[\Ker(f) \hookrightarrow M] \to \Ker[N \twoheadrightarrow \Coker(f)] = \Image(f) \]
	merupakan isomorfisme, tetapi hal ini cukup diperiksa pada tingkat
	$\cated{Mod}B$.

	Berdasarkan konstruksi, $U$ mempertahankan kernel dan kokernel, sehingga
	eksak; jelas pula bahwa $U$ setia.
\end{proof}

% segment-id: o014.aljabr2.chapter7.morita-theory-b.p009
Kembali ke pokok pembahasan, kita lanjutkan pengkajian funktor
$(\mathord\cdot)\dotimes{A}P$.

% Reference: Deligne, Categories tannakiennes, 4.5 Proposition.
% segment-id: o014.aljabr2.chapter7.morita-theory-b.pr003
\begin{proposition}\label{prop:Morita-comonad-ft}
	Misalkan bimodul-$(A,B)$ $P$, sebagai modul-$B$ kanan, projektif dan
	dibangkitkan secara hingga, dan misalkan $(\mathord\cdot)\dotimes{A}P$
	merupakan funktor eksak setia. Maka modul-$A$ kanan $N$ dibangkitkan secara
	hingga jika dan hanya jika modul-$B$ kanan $N\dotimes{A}P$ dibangkitkan secara
	hingga.
\end{proposition}
% segment-id: o014.aljabr2.chapter7.morita-theory-b.pf004
\begin{proof}
	Untuk arah ``hanya jika'', andaikan terdapat $n\in\Z_{\geq0}$ dan
	homomorfisme surjektif $A^{\oplus n}\twoheadrightarrow N$. Setelah
	menerapkan funktor, kita memperoleh homomorfisme surjektif
	$P^{\oplus n}\twoheadrightarrow N\dotimes{A}P$, sedangkan $P$, sebagai
	modul-$B$ kanan, dibangkitkan secara hingga.

	Untuk arah ``jika'', ambil submodul $N'$ dari $N$ yang dibangkitkan secara
	hingga sedemikian sehingga citra $N'\dotimes{A}P$ memuat suatu himpunan
	pembangun bagi $N\dotimes{A}P$ sebagai modul-$B$. Karena
	$(\mathord\cdot)\dotimes{A}P$ eksak setia, sesungguhnya kita mempunyai
	$N'\dotimes{A}P\rightiso N\dotimes{A}P$, dan selanjutnya $N'=N$.
\end{proof}

% segment-id: o014.aljabr2.chapter7.morita-theory-b.p010
Sesungguhnya, sifat bahwa suatu modul dibangkitkan secara hingga dapat dicirikan
sepenuhnya dalam bahasa kategori sebagai berikut.
